\section{Introduction}

The foundational role of group theory in Quantum Field Theory (QFT) was
established by Wigner \cite{Wigner1939}, who showed that physical states are
defined as unitary irreducible representations of the inhomogeneous Lorentz
group (Poincar\'{e} group). This classification, based on the ``little group''
method, reveals that fundamental wave equations emerge as specific realizations
of the underlying symmetry structure: scalar fields (spin 0) are governed by the
Klein-Gordon equation, while vector fields (spin 1) correspond to the Maxwell
and Proca equations \cite{ryderQuantumFieldTheory1996}. Within this framework, the Dirac equation
occupies a role for describing fermionic fields; while the modern
group-theoretical approach derives it directly from the transformation
properties of spinors under the Lorentz group, its original formulation was
motivated by the necessity to resolve ``duplexity'' phenomena in atomic spectra,
leading Dirac to construct a wave equation linear in time and momentum that
satisfied the requirements of both relativity and quantum mechanics
\cite{diracQuantumTheoryElectron1928}.

In recent developments, Watson and Musielak \cite{watsonChiralSymmetryDirac2020,
watsonChiralDiracEquation2021, Watson2022} revisited the
group-theoretical foundations of the field description. Instead of relying on
the standard factorization of the Klein-Gordon equation, they derived the field
equations directly from the irreducible representations of the Poincar\'{e}
group extended by parity. This approach revealed an intrinsic freedom in the
choice of the chiral basis, parametrized by a new physical degree of freedom
identified as the chiral angle. By identifying this degree of freedom, they
obtained a generalized form of the Dirac equation, hereafter referred to as the
Chiral Dirac Equation (CDE), which incorporates a pseudoscalar mass term
alongside the standard scalar mass. The authors suggested that this additional
chiral degree of freedom has physical implications; specifically, the interplay
between the scalar and pseudoscalar mass components offers a mechanism to
explain the suppression of neutrino masses, proposing the new massive particle
as a candidate for Dark Matter.

The first critical issue concerns the consistency of the quantum mechanical
description derived from the CDE. Standard quantum mechanics postulates that
physical observables must be represented by self-adjoint operators to guarantee
real energy eigenvalues and unitary time evolution (probability conservation).
While the Dirac equation is consistent with these principles, the introduction
of a chiral angle and a pseudoscalar mass term modifies the structure of the
evolution operator. It is not immediately evident that the Hamiltonian
associated with this generalized equation retains hermiticity under the new
chiral basis transformations. Therefore, we aim to explicitly derive the
Hamiltonian from the CDE and test its hermiticity to ensure the theory's
physical viability.

Secondly, a significant algebraic inconsistency arises when applying the
framework to neutrino physics under strict chirality constraints \cite{watsonChiralSymmetryDirac2020}.
Preliminary analyses indicate that for a pure left-handed field, defined by the chiral
projection $\nu_L = \frac{1}{2}\left(1-\gamma^5\right)\nu$, the scalar bilinear
$\overline{\nu}\nu$ vanishes identically. Consequently, the Yukawa Lagrangian
$\mathcal{L}_Y$ derived in the original work becomes trivial
($\mathcal{L}_Y = 0$), thereby nullifying the proposed mechanism for generating
massive neutrinos in the Dirac context. To address this, we propose to explore
the Majorana mass term within the context of a more general see-saw mechanism,
functioning as a chiral see-saw.

Finally, while the original derivation focused on bispinors, effectively
restricting the analysis to spin-1/2 particles
\cite{watsonChiralSymmetryDirac2020}, Watson et al. have also explored
vector bosons \cite{Watson2022}. However, the method employed---based on the irreducible
representations of the Poincar\'{e} group---is fundamentally geometric and
should be applicable to other spin sectors. To this end, we propose a
review of the feasibility of extending this chiral formalism to
higher-spin fields. Specifically, we examine whether similar chiral degrees of
freedom can be consistently introduced into the Proca equation for massive
vector bosons (spin-1) or the Rarita-Schwinger fields (spin-3/2).
\section{Theoretical Framework}
\label{sec:Framework}

We adopt the formalism of relativistic quantum field theory on a flat,
non-interacting Minkowski spacetime. The background geometry is defined by the
orthonormal basis $\{ \gamma^\mu \}$ satisfying the Clifford algebra
$\mathcal{C}\ell_{1,3}(\mathbb{C})$, defined by the fundamental anticommutation
relation
\begin{equation}
	\{ \gamma^\mu, \gamma^\nu \} \equiv \gamma^\mu \gamma^\nu + \gamma^\nu
	\gamma^\mu = 2 \eta^{\mu \nu},
\end{equation}
where $\eta^{\mu \nu} = \text{diag}(1, -1, -1, -1)$ is the Minkowski metric.
Consequently, $(\gamma^0)^2 = 1$ and $(\gamma^k)^2 = -1$. Greek indices $\mu =
0, \dots, 3$ denote spacetime coordinates, while Latin indices $k = 1, \dots, 3$
denote spatial components.

The chiral structure of the theory is determined by the pseudoscalar element
$\gamma^5$, defined as
\begin{equation}
	\gamma^5 \equiv i \gamma^0 \gamma^1 \gamma^2 \gamma^3.
\end{equation}
Since the spacetime dimension is even, $\gamma^5$ anticommutes with all vector
generators $\gamma^\mu$ and squares to the identity, $(\gamma^5)^2 = 1$. This
algebraic property permits the construction of the orthogonal chiral projectors
\begin{equation}
	P_L = \frac{1}{2}(1 - \gamma^5), \quad P_R = \frac{1}{2} (1 + \gamma^5),
\end{equation}
which satisfy the idempotency relations $P_{L,R}^2 = P_{L,R}$ and the
orthogonality condition $P_L P_R = 0$. A spinor field $\psi$ therefore
decomposes uniquely into left- and right-handed chiral components, $\psi = P_L
\psi + P_R \psi \equiv \psi_L + \psi_R$.

\subsection{The Chiral Dirac Ansatz}

Standard formulations of the Dirac equation assume a scalar mass term that
couples chiral sectors trivially. Following Watson and Musielak
\cite{watsonChiralSymmetryDirac2020}, we consider a generalized dynamical
equation invariant under the irreducible representations of the Poincar\'{e}
group extended by parity. This yields the Chiral Dirac Equation (CDE) for a
spinor field $\psi$
\begin{equation}
	(i \gamma^\mu \partial_\mu - m e^{- i \alpha \gamma^5}) \psi = 0.
	\label{eq:CDE}
\end{equation}
Here, $m$ represents the rest mass, and $\alpha$ is a chiral angle parametrizing
the rotation between the scalar and pseudoscalar mass sectors. The corresponding
Lorentz-invariant effective Lagrangian density is given by
\begin{equation}
	\mathcal{L}_{\text{eff}} = \bar{\psi} (i \gamma^\mu \partial_\mu - m
	e^{- i \alpha \gamma^5}) \psi,
	\label{eq:Lagrangian}
\end{equation}
where $\bar{\psi} \equiv \psi^\dagger \gamma^0$ is the Dirac adjoint. This
Lagrangian serves as the generating functional for the dynamics and
conservation laws investigated in the subsequent sections.
\section{Hamiltonian Consistency and Hermiticity}
\label{sec:Hermiticity}

To ensure the physical viability of the Chiral Dirac Equation (CDE), we must verify
that it generates a unitary time evolution. This requirement demands that the
Hamiltonian operator governing the dynamics be self-adjoint on the Hilbert space
of states. We proceed by explicitly deriving the Hamiltonian from the effective
Lagrangian density and analyzing its properties under Hermitian conjugation.

\subsection{Canonical Derivation}

The canonical momentum conjugate to the field $\psi$ is obtained from the
Lagrangian density $\mathcal{L}_{\text{eff}}$ (Eq. \ref{eq:Lagrangian}) via
differentiation with respect to the time derivative of the field, $\dot{\psi}
\equiv \partial_0 \psi$
\begin{equation}
	\pi = \frac{\partial \mathcal{L}_{\text{eff}}}{\partial \dot{\psi}} =
	i \psi^\dagger.
\end{equation}
The Hamiltonian density $\mathcal{H}$ is constructed via the Legendre transform
\begin{equation}
	\mathcal{H} = \pi \dot{\psi} - \mathcal{L}_{\text{eff}}.
\end{equation}
Substituting the explicit form of $\mathcal{L}_{\text{eff}}$ and isolating the
temporal component of the derivative term, we find
\begin{align}
	\mathcal{H} &= i \psi^\dagger \dot{\psi} - \psi^\dagger \gamma^0
	(i \gamma^0 \dot{\psi} + i \gamma^k \partial_k \psi - m e^{-i \alpha \gamma^5} \psi)
	\nonumber \\
							&= i \psi^\dagger \dot{\psi} - i \psi^\dagger \dot{\psi} - i \psi^\dagger \gamma^0
							\gamma^k \partial_k \psi + \psi^\dagger \gamma^0 m e^{-i \alpha \gamma^5} \psi
							\nonumber \\
							&= \psi^\dagger \left( -i \gamma^{0k} \partial_k + \gamma^0 m e^{-i \alpha \gamma^5}
							\right) \psi,
\end{align}
where we have introduced the Hermitian spatial generators $\gamma^{0k} \equiv
\gamma^0 \gamma^k$. The quantum mechanical Hamiltonian operator $H$, defined by
the Schrödinger-like evolution equation $i \partial_0 \psi = H \psi$, is
therefore identified as
\begin{equation}
	H = \gamma^{0k} p_k + \gamma^0 m e^{-i \alpha \gamma^5},
	\label{eq:HamiltonianOperator}
\end{equation}
where $p_k = -i \partial_k$ is the spatial momentum operator.

\subsection{Proof of Hermiticity}

For the theory to be consistent, $H$ must be Hermitian ($H = H^\dagger$). We
examine the adjoint of the derived operator
\begin{equation}
	H^\dagger = \left( \gamma^{0k} p_k \right)^\dagger + \left( \gamma^0 m
	e^{-i \alpha \gamma^5} \right)^\dagger.
\end{equation}
The first term is the standard kinetic energy operator. Since the momentum $p_k$
is Hermitian and the matrix $\gamma^{0k}$ is Hermitian, as $(\gamma^0 \gamma^k)^\dagger
= \gamma^{k\dagger} \gamma^{0\dagger} = (-\gamma^k)(\gamma^0) = \gamma^0
\gamma^k$,
their product is Hermitian.

The second term, containing the chiral mass correction, requires careful
algebraic treatment. Expanding the adjoint yields
\begin{equation}
	\left( \gamma^0 m e^{-i \alpha \gamma^5} \right)^\dagger = m
	\left( e^{-i \alpha \gamma^5} \right)^\dagger \left( \gamma^0 \right)^\dagger.
\end{equation}
Utilizing the hermiticity of $\gamma^5$ and $\gamma^0$, and the identity
$(e^A)^\dagger = e^{A^\dagger}$, we obtain
\begin{equation}
	m e^{i \alpha \gamma^5} \gamma^0.
\end{equation}
Critically, $\gamma^0$ and $\gamma^5$ anticommute ($\{ \gamma^0, \gamma^5 \} =
0$). This anticommutation implies the commutation relation for the exponential
function
\begin{equation}
	e^{i \theta \gamma^5} \gamma^0
	= \left(\cos{\theta} + i \gamma^5 \sin{\theta}\right) \gamma^0
	= \gamma^0 \left(\cos{\theta} - i \gamma^5 \sin{\theta}\right)
	= \gamma^0 e^{-i \theta \gamma^5}.
\end{equation}
Applying this commutation to the mass term, we recover the original operator
\begin{equation}
	m e^{i \alpha \gamma^5} \gamma^0 = \gamma^0 m e^{-i \alpha \gamma^5}.
\end{equation}
Thus,
\begin{equation}
	H^\dagger = \gamma^{0k} p_k + \gamma^0 m e^{-i \alpha \gamma^5} = H.
\end{equation}
We conclude that the Hamiltonian of the Chiral Dirac Equation is strictly
Hermitian. The introduction of the chiral angle $\alpha$ preserves the
self-adjoint nature of the energy operator, ensuring real energy eigenvalues and
unitary time evolution without requiring a modification of the inner product or
metric.
\section{Chiral Seesaw Mechanism}

As an alternative approach to explain the smallness of neutrino masses, we
revisit the mass generation mechanism proposed by Watson and Musielak within the
framework of the Chiral Dirac Equation (CDE)
\cite{watsonChiralSymmetryDirac2020, watsonChiralDiracEquation2021}.
In this formalism, the fermion mass arises from Yukawa couplings to both a
scalar field $\phi_1$ (Standard Model-like) and a pseudoscalar field $\phi_2$.
In Watson's proposal \cite{Watson2022}, restricting the theory to pure
left-handed fermionic fields causes the mass terms to vanish identically.
Considering the original Yukawa Lagrangian postulated by Watson
\begin{equation} \label{eq:original_yukawa}
  \mathcal{L}_Y = -\frac{\lambda_1 v_1}{\sqrt{2}}\,\overline{\nu}\nu -
  \frac{\lambda_2 v_2}{\sqrt{2}}\,\overline{\nu}\gamma^{5}\nu.
\end{equation}

Both terms vanish if we impose the left-handed chirality constraint on the
neutrino field ($P_L \nu = \nu$, $P_R \nu = 0$). The scalar product
$\overline{\nu}\nu$ vanishes because the Dirac adjoint of a left-handed field
behaves as a right-handed projector acting to the left
\begin{equation}
\begin{aligned}
  \overline{\nu} &= \overline{P_L \nu} \\
                 &= \nu^\dagger P_L \gamma^0 \\
                 &= \nu^\dagger \gamma^0 P_R \\
                 &= \overline{\nu}\, P_R .
\end{aligned}
\end{equation}

Multiplying by $\nu$ from the right, and recalling that $\nu = P_L \nu$, we
obtain
\begin{equation}
  \overline{\nu}\, \nu = \left(\overline{\nu}\, P_R\right) \left(P_L
  \nu\right) = \overline{\nu} \underbrace{\left(P_R P_L\right)}_{0} \nu = 0.
\end{equation}

Similarly, for the pseudoscalar term $\overline{\nu} \gamma^{5} \nu$, we utilize
the fact that $\gamma^5$ commutes with the chiral projectors.
Following the previous result
\begin{equation}
\begin{aligned}
  \overline{\nu} \gamma^{5} \nu &= \left(\overline{\nu} P_R\right) \gamma^{5}
  \left(P_L \nu\right) \\
                                &= \overline{\nu} \gamma^{5} \left(P_R P_L\right)
                                \nu \\
                                &= 0.
\end{aligned}
\end{equation}

Thus, without the inclusion of a right-handed component, no Dirac mass term can
be generated in this framework. However, if we assume the existence of a
right-handed neutrino component $\nu_R$ (such that $\nu = \nu_L + \nu_R$), the
Dirac mass term does not vanish. Instead, it acquires a complex phase determined
by the chiral angle $\alpha$, which is intrinsic to the CDE. To naturally
suppress the neutrino mass without relying on fine-tuning, we invoke the
Majorana mass term, leading to a Chiral Seesaw Mechanism \cite{peskinIntroductionQuantumField2019}.


\subsection{The Chiral Dirac Mass Term}

Following Watson and Musielak \cite{watsonChiralSymmetryDirac2020},
the general Yukawa Lagrangian expanded about the vacuum expectation values
(VEVs) $v_1$ and $v_2$ corresponds to \Cref{eq:original_yukawa}. Defining the
scalar mass parameter $M = \lambda_1 v_1 / \sqrt{2}$ and the pseudoscalar mass
parameter $\tilde{M} = \lambda_2 v_2 / \sqrt{2}$, and decomposing the fields
into their chiral components, the Dirac mass term
becomes
\begin{equation}
  \mathcal{L}_{\text{Dirac}} = \overline{\nu}_L \left(M + \tilde{M}\right)
  \nu_R + \overline{\nu}_R \left(M - \tilde{M}\right) \nu_L.
\end{equation}

In the CDE formalism, these parameters are related to the propagation mass $m$
and the chiral angle $\alpha$ by $M = m \cos \alpha$ and $\tilde{M} = -im \sin
\alpha$. Consequently, the effective Dirac mass coupling $M_D$ connecting
$\overline{\nu}_L$ and $\nu_R$ is identified as
\begin{equation} \label{eq:chiral_M_D}
  M_D \equiv M + \tilde{M} = m\left(\cos \alpha - i \sin \alpha\right) = m
  e^{-i\alpha}.
\end{equation}

The result in \eqref{eq:chiral_M_D} is crucial, the Dirac mass in this model
carries a physical phase derived from the underlying chiral symmetry breaking.

\subsection{The Full Majorana-Dirac Lagrangian}

To implement the seesaw mechanism, we introduce Majorana mass terms. Unlike
Dirac terms which couple left- and right-handed fields ($\overline{\nu}_L
\nu_R$), Majorana mass terms violate lepton number by two units. Specifically,
we utilize the charge-conjugated right-handed field $\nu_R^c =
C\overline{\nu}_R^T$. The most general Lorentz-invariant mass Lagrangian
including both the Chiral Dirac term derived above and the Majorana terms is
\begin{equation} \label{eq:L_mass_full}
  \mathcal{L}_{\text{Mass}}
  =
  \overline{\nu}_L M_D \nu_R
  +
  \frac{1}{2}\overline{\nu_L^c} M_L \nu_L
  +
  \frac{1}{2}\overline{\nu_R^c} M_R \nu_R
  +\text{h.c.}
\end{equation}

Here, $M_L$ represents the Majorana mass for the active left-handed neutrinos
(assumed to be negligible or zero in the minimal Type-I seesaw), and $M_R$ is
the heavy Majorana mass scale for the sterile neutrinos.

\subsection{The Chiral Mass Matrix}

To find the physical mass eigenstates, we rewrite \eqref{eq:L_mass_full} using
the basis of left-handed fields. We define the interaction basis spinor $\Psi =
\left(\nu_L, \nu_R^c\right)^T$. Utilizing the Majorana identity, the Lagrangian
allows the explicit matrix expansion
\begin{equation}
  \mathcal{L}_{\text{Mass}} = \frac{1}{2} \overline{\Psi^c} \mathcal{M} \Psi +
  \text{h.c.},
\end{equation}

where the mass matrix $\mathcal{M}$ incorporates the chiral modifications from
the CDE defined in \eqref{eq:chiral_M_D}
\begin{equation} \label{eq:mass_matrix}
  \mathcal{M} =
  \begin{pmatrix}
    M_L & M_D \\
    M_D^T & M_R
  \end{pmatrix}
  =
  \begin{pmatrix}
    M_L & m e^{-i\alpha} \\
    m e^{-i\alpha} & M_R
  \end{pmatrix}.
\end{equation}

\subsection{The Neutrino Chiral Mass}

To account for the smallness of neutrino masses, we adopt the following
identifications in the matrix \eqref{eq:mass_matrix}
\begin{enumerate}
  \item \textbf{$M_R$:} Heavy sterile scale ($M_R \gg m$).
  \item \textbf{$M_L$:} Assumed $M_L \approx 0$.
  \item \textbf{$M_D$:} Given by $m e^{-i\alpha}$.
\end{enumerate}

Based on these assumptions, the mass matrix is given by
\begin{equation} \label{eq:mass_matrix_approx}
  \mathcal{M} \approx \begin{pmatrix} 0 & m e^{-i\alpha} \\ m e^{-i\alpha} &
  M_R \end{pmatrix}.
\end{equation}

Finding the eigenvalues of \eqref{eq:mass_matrix_approx} and applying the seesaw
approximation ($m \ll M_R$), we obtain two mass terms
\begin{enumerate}
  \item \textbf{Heavy Mass ($m'_1$):} Sterile Dominant State.
    \begin{equation}
      m'_1 \approx M_R + \frac{m^2 e^{-2i\alpha}}{M_R} \approx M_R
    \end{equation}

  \item \textbf{Light Mass ($m'_2$):} Active Dominant State.
    \begin{equation} \label{eq:light_mass_eigenvalue}
      m'_2 \approx - \frac{\left(m e^{-i\alpha}\right)^2}{M_R} = -
      \frac{m^2}{M_R} e^{-2i\alpha}
    \end{equation}
\end{enumerate}

\subsection{Mass Suppression and CP Violation}

The application of the Chiral Seesaw Mechanism yields the following expression
for the active neutrino mass based on \eqref{eq:light_mass_eigenvalue}
\begin{equation} \label{eq:final_neutrino_mass}
  m'_{\nu} \approx    \frac{m^2}{M_R} e^{-2i\alpha}.
\end{equation}

This equation implies two significant physical results.

\subsubsection{Mass Suppression}

The magnitude of the active neutrino mass is given by $\left|m'_{\nu}\right|
\approx m^2 / M_R$. This confirms the standard seesaw suppression mechanism. For
instance, if we assume $m \sim \qty{100}{\giga\electronvolt}$ and
$\left|m'_{\nu}\right| \sim \qty{0.1}{\electronvolt}$, the mechanism
necessitates a heavy scale of $M_R \sim \qty{e14}{\giga\electronvolt}$.

\subsubsection{The Chiral Phase and CP Violation}

A novel feature of this model is that the phase of the light neutrino
mass in \eqref{eq:final_neutrino_mass} is directly determined by the chiral
angle $\alpha$
\begin{equation}
  \text{Phase}\left(m'_{\nu}\right) = e^{-2i\alpha}.
\end{equation}

This dependence has the following implications:
\begin{itemize}
  \item \textbf{Leptonic CP Violation:} In a three-generation scenario, this
    complex phase translates into Majorana phases in the PMNS mixing matrix. A
    value of $\alpha \neq k\pi/2$ implies a novel source of leptonic CP
    violation.
  \item \textbf{BSM Link:} This establishes a direct link between the
    theoretical angle $\alpha$ of the CDE and a physical observable mediated by
    the seesaw scale $M_R$.
\end{itemize}
\section{Generalization to Higher Spins}
\label{sec:HigherSpins}

To systematically address the extension to higher spins, we adopt the
Bargmann-Wigner (BW) formalism, which provides a rigorous group-theoretical
framework for constructing relativistic wave equations for particles of
arbitrary spin $s$. Established by V. Bargmann and E. P. Wigner
\cite{bargmannGroupTheoreticalDiscussion1997}, this method posits that the wavefunction of a particle
with spin $s$ can be represented by a multispinor field $\Psi_{\alpha_1 \alpha_2
\dots \alpha_{2s}}$ of rank $2s$, which effectively describes a composite state
of $2s$ fundamental spin-1/2 fermions. To ensure that this object transforms
irreducibly under the Poincar\'{e} group, the multispinor must satisfy two
conditions: it must be totally symmetric under the exchange of any pair of
spinor indices, and it must satisfy the Dirac equation for each index
independently
\begin{equation}
  \left( i\,\gamma^\mu \partial_\mu - m \right)_{\alpha_k \beta}\,
  \Psi_{\alpha_1 \dots \beta \dots \alpha_{2s}}(x)
  = 0,
  \qquad \text{for } k = 1, \dots, 2s .
\end{equation}

To address the specific extension to vector bosons, we adopt the methodology
established by Watson in their generalization of the Proca equation
\cite{Watson2022}. Their approach serves as the guiding framework for this
analysis. Specifically, for the vector boson case ($s=1$), the formalism begins
with the Bargmann-Wigner equations for a symmetric rank-2 spinor
$\Psi_{\alpha\beta}$. Crucially, the standard scalar mass $m$ is replaced by the
chiral mass operator $\hat{M}$ derived from the fermionic theory, applied to
each spinor index independently
\begin{equation}
  \left(i \gamma^\mu \partial_\mu - \hat{M}\right)_{\alpha \lambda}
  \Psi_{\lambda \beta}(x) = 0, \quad \text{with} \quad \hat{M} = m
  e^{-i\alpha\gamma^5}.
\end{equation}

The wavefunction is then expanded in terms of the Clifford algebra generators.
The symmetry constraint intrinsic to the Bargmann-Wigner formalism for $s=1$
requires the spinor $\Psi_{\alpha\beta}$ to be symmetric, which effectively
eliminates the scalar, pseudoscalar, and axial-vector components. Consequently,
the wavefunction reduces naturally to a superposition of only the vector
($A_\mu$) and tensor ($F_{\mu\nu}$) sectors
\begin{equation}
  \Psi(x) = \left( \gamma^\mu A_\mu(x) + \frac{1}{2}\sigma^{\mu\nu}
  F_{\mu\nu}(x) \right) C,
\end{equation}
where $C$ is the charge conjugation matrix. Substituting this ansatz back into
the fundamental chiral equations leads to a decoupling of the components. While
the standard theory yields the Proca equation $\partial_\mu F^{\mu\nu} + m^2
A^\nu = 0$, the chiral extension results in the Generalized Proca Equation
\begin{equation}
  \partial_\mu F^{\mu\nu} + m^2 \cos^2\left(\alpha\right) A^\nu + \dots = 0.
\end{equation}

It is important to clarify that the appearance of the $\cos^2\left(\alpha\right)$
factor in this sector does not violate the relativistic structure of the theory;
rather, it acts as a scaling parameter that defines an effective physical mass
$M_{\text{eff}} = m \cos\left(\alpha\right)$. This implies that the chiral angle
can modulate the observable mass of the boson without altering the underlying
geometric consistency of the Poincar\'{e} group.

\subsection{The Chiral Rarita-Schwinger Equation}

Building on this success, we apply the procedure to the Rarita-Schwinger field
($s=3/2$). To maintain full consistency with the underlying Bargmann-Wigner
theory, we treat the field strictly as a totally symmetric rank-3 spinor
$\Psi_{\alpha\beta\gamma}$. In this formalism, the field is a composite state of
three fundamental spin-1/2 fermions. The chiral mass generation mechanism is
applied to each of the three constituent spinor indices. Since the chiral
rotation $e^{-i\alpha\gamma^5}$ induces a projection factor of
$\cos\left(\alpha\right)$ for a single spinor (as seen in the Dirac case), the
composite mass term for a rank-3 object must scale as the product of the
projections of its constituents. This leads to a cubic scaling law
\begin{equation}
  \hat{M}_{\text{composite}} \rightarrow m \cos^3\left(\alpha\right).
\end{equation}

Consequently, the modified Chiral Rarita-Schwinger equation takes the form
\begin{equation}
  \epsilon^{\mu\nu\rho\sigma} \gamma^5 \gamma_\nu \partial_\rho \psi_\sigma +
  i m \cos^3\left(\alpha\right) \sigma^{\mu\nu} \psi_\nu = 0.
\end{equation}

This result suggests a general scaling law for the effective mass of a particle
with spin $s$: $M_{\text{eff}}^{(s)} \approx m \cos^{2s}\left(\alpha\right)$.
This theoretical prediction aligns with the linear scaling found for Dirac
fermions ($2s=1 \rightarrow \cos^1$) and the quadratic scaling found for Proca
bosons ($2s=2 \rightarrow \cos^2$), offering a unified geometric interpretation
of mass modulation across the spin spectrum. Finally, we note that for
higher-spin fields, the introduction of non-trivial mass matrices requires
careful checking of the constraint algebra ($\gamma^\mu\psi_\mu=0$ and
$\partial^\mu\psi_\mu=0$) to avoid the propagation of acausal modes. This is a
known issue in massive Rarita-Schwinger theories referred to as the
Velo-Zwanziger problem \cite{napsucialeSpinFrac32$RaritaSchwinger2006}.
While a full causality analysis is beyond the scope of this proposal, future
work must verify the preservation of these constraints under chiral mass
modifications.

\subsection{Discussion on Generalization via Projection Operators}

The explicit construction of equations for higher spins (like $s=3/2$ or $s=2$)
using standard component notation becomes exponentially complex. To solve this,
Watson and Musielak propose a powerful generalization based on Orthogonal
Idempotents. Instead of manually constructing the Lagrangians, this method
utilizes projection operators $\hat{P}$ to define the physical subspaces
directly. For a system of $N$ constituent spinors (which would form a spin-$N/2$
particle), the generalized projection operator is defined as the tensor product
of the individual helicity projectors
\begin{equation}
  \hat{P}_{s_1, s_2, \dots, s_N}\left(\hat{n}\right) =
  \hat{P}_{s_1}\left(\hat{n}\right) \otimes \hat{P}_{s_2}\left(\hat{n}\right)
  \otimes \dots \otimes \hat{P}_{s_N}\left(\hat{n}\right),
\end{equation}
where each $\hat{P}_{s_i}\left(\hat{n}\right)$ projects the spin state relative
to a quantization axis $\hat{n}$. For the specific case of the Rarita-Schwinger
field ($N=3$), one selects the symmetric subspace of three entangled spinors.
This projector formalism is superior because it automatically handles the chiral
constraints. By defining the chiral basis for each projector (parameterized by
angles $\alpha_1, \alpha_2, \dots$), the resulting equation of motion emerges
naturally with the mass terms already modulated by the relative chiral angles.
Thus, the effective mass $M_{\text{eff}}$ is not an ad-hoc insertion but an
eigenvalue of the generalized projection operator acting on the vacuum state
\begin{equation}
  \hat{P}_{\text{total}} \Psi = M_{\text{eff}} \Psi.
\end{equation}
This confirms that the chiral degrees of freedom are fundamental to the
algebraic structure of the Poincar\'{e} group representations, providing a
path to derive wave equations for any arbitrary spin $s$.
